\documentclass{beamer}

\usetheme{Berlin}
\usecolortheme{Orchid}
\useoutertheme{miniframes}
\setbeamertemplate{navigation symbols}{}

\usepackage[utf8]{inputenc}
\usepackage[T1]{fontenc}
\usepackage{amsmath}
\usepackage{amssymb}
\usepackage{booktabs}
\usepackage{graphicx}
\graphicspath{{../images/}}

\usepackage{xcolor}
\usepackage{listings}
\usepackage{hyperref}
\usepackage{tikz}
\usetikzlibrary{arrows.meta,positioning,shapes.geometric,calc}

\lstset{
  language=Java,
  basicstyle=\ttfamily\small,
  keywordstyle=\color{blue},
  commentstyle=\color{gray},
  stringstyle=\color{teal},
  showstringspaces=false,
  columns=fullflexible,
  breaklines=true
}

% Per-slide source line (references shown on the slide itself).
\newcommand{\slidesources}[1]{%
  \vfill
  {\tiny\color{gray}\raggedright\sloppy\parbox{\linewidth}{Sources: #1}\par}%
}

% Unit-wise source strings for this subject (used by slides/notes).
% Path is relative to the lecture `latex/` folder (the compilation CWD).
% OOPS with Java (CSEG1044) — unit-wise sources
%
% These are short, student-facing reference strings intended to be used with:
%   - \slidesources{...}  (in beamer slides)
%   - \notesources{...}   (in lecture notes)
%
% Style inspiration: other subjects in My-Subjects/ use semicolon-separated sources
% and include an "accessed" date for web documentation.

\newcommand{\OOPSUnitISources}{Schildt, \textit{Java: A Beginner's Guide} (10e), McGraw-Hill (Java fundamentals + OOP basics).; Oracle Java Tutorials: Getting Started + Language Basics + Classes and Objects (accessed \today).; Java SE API docs (e.g., \texttt{String}, \texttt{Scanner}) (accessed \today).}

\newcommand{\OOPSUnitIISources}{Schildt, \textit{Java: The Complete Reference} (12e), McGraw-Hill (classes, inheritance, packages, interfaces).; Bloch, \textit{Effective Java} (3e), Addison-Wesley, 2018 (design choices; interfaces vs abstract classes).; Oracle Java Tutorials: Classes and Objects + Interfaces + Packages (accessed \today).; Java SE API docs (e.g., \texttt{Comparable}, \texttt{Runnable}) (accessed \today).}

\newcommand{\OOPSUnitIIISources}{Schildt, \textit{Java: The Complete Reference} (12e) (nested/inner classes, exceptions, threads, I/O).; Oracle Java Tutorials: Nested Classes + Exceptions + Concurrency + Basic I/O (accessed \today).; Java SE API docs (e.g., \texttt{Thread}, \texttt{AutoCloseable}) (accessed \today).}

\newcommand{\OOPSUnitIVSources}{Schildt, \textit{Java: The Complete Reference} (12e) (generics, lambdas, Swing, JDBC).; Bloch, \textit{Effective Java} (3e) (generics + lambdas).; Oracle Java Tutorials: Generics + Lambda Expressions + Swing + JDBC Basics (accessed \today).; Java SE API docs (e.g., \texttt{java.util.function}, \texttt{javax.swing}, \texttt{java.sql}) (accessed \today).}

\newcommand{\OOPSUnitVSources}{Schildt, \textit{Java: The Complete Reference} (12e) (Collections Framework + wrapper classes).; Oracle Java docs: Collections Framework overview (accessed \today).; Java SE API docs (\texttt{java.util} collections; wrapper classes in \texttt{java.lang}) (accessed \today).}

\newcommand{\OOPSUnitVISources}{Git SCM: \textit{Pro Git} (2e) (version control workflow), accessed \today.; GitHub Docs (repositories, issues, pull requests), accessed \today.; Oracle Java Tutorials: Swing + JDBC (accessed \today).; JUnit 5 User Guide (accessed \today).; Apache Maven Project: Getting Started (accessed \today).; Gradle User Manual: Getting Started (accessed \today).; UML class diagrams: UML 2.x overview/spec (accessed \today).}

\newcommand{\OOPSMiscWhyInterfacesSources}{Schildt, \textit{Java: The Complete Reference} (12e) (interfaces).; Bloch, \textit{Effective Java} (3e) (prefer interfaces where appropriate).; Oracle Java Tutorials: Interfaces (accessed \today).; Java SE API docs (e.g., \texttt{Comparable}, \texttt{Runnable}) (accessed \today).}



\title[OOPS with Java]{Object Oriented Programming (CSEG1044)}
\subtitle{Unit 02 -- Lecture 04: Inheritance, `super`, Overriding, and Dynamic Dispatch}
\begin{document}

\begin{frame}[fragile]
  \titlepage
  \vspace{0.5em}
  \slidesources{\OOPSUnitIISources}
\end{frame}

\begin{frame}{Quick Links}
  \centering
  \hyperlink{sec:overview}{\beamerbutton{Overview}}\hspace{0.6em}
  \hyperlink{sec:concepts}{\beamerbutton{Key Topics}}\hspace{0.6em}
  \hyperlink{sec:demo}{\beamerbutton{Demo}}\hspace{0.6em}
  \hyperlink{sec:summary}{\beamerbutton{Summary}}
  \slidesources{\OOPSUnitIISources}
\end{frame}

\begin{frame}{Agenda}
  \tableofcontents
  \slidesources{\OOPSUnitIISources}
\end{frame}

\section{Overview}
\label{sec:overview}

\begin{frame}{Lecture Snapshot}
\begin{itemize}[<+->]
  \item \textbf{Unit focus:} Classes, Inheritance, Packages, and Interfaces
  \item \textbf{Course thread:} Use Inheritance, `super`, Overriding, and Dynamic Dispatch to strengthen the domain model, APIs, and access rules inside Campus Resource Manager.
\end{itemize}
  \slidesources{\OOPSUnitIISources}
\end{frame}

\begin{frame}{Recall Before You Start}
\begin{itemize}[<+->]
  \item \textbf{Class design}: Inheritance helps only when the parent and child already have meaningful responsibilities. Main reference: \texttt{Unit 02 / Lecture 01 slides/notes}.
  \item \textbf{Constructors}: The rule for super(...) depends on the constructor logic from the previous lecture. Main reference: \texttt{Unit 02 / Lecture 02 slides/notes}.
\end{itemize}
  \slidesources{\OOPSUnitIISources}
\end{frame}


\begin{frame}{Learning Outcomes}
  \begin{itemize}[<+->]
    \item Decide is-a (inheritance) vs has-a (composition).
    \item Use \texttt{super(...)} and \texttt{super.method()} correctly.
    \item Override methods and explain dynamic dispatch with an example.
  \end{itemize}
  \slidesources{\OOPSUnitIISources}
\end{frame}

\begin{frame}{Key Terms to Know}
\small
\begin{columns}[t]
\column{0.48\textwidth}
\begin{itemize}
  \item \textbf{Inheritance}: a subclass extends a superclass (\texttt{extends}).
  \item \textbf{Superclass/Subclass}: parent/child in inheritance hierarchy.
  \item \textbf{Overriding}: subclass provides its own implementation of a superclass method (same signature).
\end{itemize}
\column{0.48\textwidth}
\begin{itemize}
  \item \textbf{Dynamic dispatch}: method call resolved at runtime based on the actual object type.
  \item \textbf{Upcasting}: treating a subclass object as a superclass reference, for example \path|Shape s = new Circle(...)|.
  \item \textbf{Composition}: has-a relationship (\texttt{Car} has an \texttt{Engine}).
\end{itemize}
\end{columns}
  \slidesources{\OOPSUnitIISources}
\end{frame}

\begin{frame}{Study Flow for This Lecture}
\begin{enumerate}[<+->]
  \item Refresh the prerequisite bridge before reading new ideas in isolation.
  \item Study the main build in this order: \textit{Inheritance and types (conceptual)} $\rightarrow$ \textit{\texttt{super} usage} $\rightarrow$ \textit{Overriding and dynamic dispatch}.
  \item Trace the worked example and connect each line back to one concept frame.
  \item Attempt the mini exercises before moving to the recap and exit check.
\end{enumerate}
  \slidesources{\OOPSUnitIISources}
\end{frame}


\section{Key Topics}
\label{sec:concepts}

\begin{frame}{Unit 02: Classes, Inheritance, Packages and Interfaces}
  \begin{itemize}[<+->]
    \item Lecture focus: Inheritance, `super`, Overriding, and Dynamic Dispatch
  \end{itemize}
  \slidesources{\OOPSUnitIISources}
\end{frame}

\begin{frame}{Today: Roadmap}
  \begin{itemize}[<+->]
    \item Inheritance and types (conceptual); "is-a" vs "has-a"
    \item `super` usage (constructors + member access)
    \item Overriding + runtime polymorphism (dynamic method dispatch)
  \end{itemize}
  \slidesources{\OOPSUnitIISources}
\end{frame}

\begin{frame}{Is-a vs Has-a (design choice)}
  \begin{itemize}[<+->]
    \item Use inheritance for \textbf{is-a}: Dog is-a Animal.
    \item Use composition for \textbf{has-a}: Car has-a Engine.
    \item Avoid inheritance only for code reuse (creates strong coupling).
  \end{itemize}
  \slidesources{\OOPSUnitIISources}
\end{frame}

\begin{frame}{\texttt{super} usage}
  \begin{itemize}[<+->]
    \item \texttt{super(...)} calls the parent constructor (must be first statement).
    \item \texttt{super.method()} calls parent implementation if overridden.
  \end{itemize}
  \slidesources{\OOPSUnitIISources}
\end{frame}

\begin{frame}{Overriding + dynamic dispatch}
  \begin{itemize}[<+->]
    \item Overriding: same signature in subclass; use \texttt{@Override}.
    \item Dynamic dispatch: method chosen at runtime based on actual object type.
    \item Example: \texttt{Shape s = new Circle(...);} then \texttt{s.area()} runs \texttt{Circle.area()}.
  \end{itemize}
  \slidesources{\OOPSUnitIISources}
\end{frame}

\begin{frame}{Checkpoint and Reflection}
\begin{block}{Reflection prompt}
Where would \textit{Inheritance, `super`, Overriding, and Dynamic Dispatch} matter most in Campus Resource Manager, and which design choice would you justify using this lecture's ideas?
\end{block}
\begin{block}{Quick self-check}
Which part needs one more example before you move on: \textit{Inheritance and types (conceptual)}, \textit{\texttt{super} usage}, the worked example, or the recap?
\end{block}
  \slidesources{\OOPSUnitIISources}
\end{frame}

\section{Demo / Example}
\label{sec:demo}

\begin{frame}[fragile,shrink=8]{Example: \texttt{Shape} + \texttt{Circle}}
\begin{lstlisting}
class Shape {
    private final String color;
    Shape(String color) { this.color = color; }
    double area() { return 0.0; }
}

class Circle extends Shape {
    private final double r;
    Circle(String color, double r) {
        super(color);
        this.r = r;
    }
    @Override double area() { return Math.PI * r * r; }
}
\end{lstlisting}
  \slidesources{\OOPSUnitIISources}
\end{frame}

\begin{frame}[fragile,shrink=8]{Dynamic dispatch demo}
\begin{lstlisting}
public class Demo {
    public static void main(String[] args) {
        Shape s = new Circle("red", 2.0);
        System.out.println(s.area());
    }
}
\end{lstlisting}
  \slidesources{\OOPSUnitIISources}
\end{frame}

\begin{frame}{Mini Exercises (2--3 mins)}
  \begin{enumerate}[<+->]
    \item Student--IDCard: is-a or has-a?
    \item What must be first in a constructor: \texttt{super(...)} or other statements?
    \item Overloading vs overriding: which is runtime?
  \end{enumerate}
  \slidesources{\OOPSUnitIISources}
\end{frame}

\begin{frame}{Watch-outs}
\begin{itemize}[<+->]
  \item Choosing inheritance just to reduce code repetition quickly.
  \item Forgetting that super(...) must run before subclass-specific initialization.
  \item Confusing overriding with overloading.
\end{itemize}
  \slidesources{\OOPSUnitIISources}
\end{frame}

\section{Summary}
\label{sec:summary}

\begin{frame}{Key Takeaways}
  \begin{itemize}[<+->]
    \item Use inheritance only for true is-a relationships.
    \item \texttt{super(...)} initializes the parent part of an object.
    \item Overriding + dynamic dispatch enable runtime polymorphism.
  \end{itemize}
  \slidesources{\OOPSUnitIISources}
\end{frame}

\begin{frame}{Checkpoint Questions}
  \begin{enumerate}[<+->]
    \item What is the difference between is-a and has-a?
    \item Why do we use \texttt{@Override}?
    \item What is dynamic dispatch?
  \end{enumerate}
  \slidesources{\OOPSUnitIISources}
\end{frame}

\begin{frame}{Next Study Steps}
\begin{itemize}[<+->]
  \item Revisit the recall references if any older idea still feels weak.
  \item Rework the lecture example without looking at the notes first.
  \item Attempt the self-study practice and then answer the exit questions in full sentences.
  \item Apply the idea once inside Campus Resource Manager to make the concept stick.
\end{itemize}
  \slidesources{\OOPSUnitIISources}
\end{frame}

\begin{frame}{Next Lecture Preview}
  \textbf{Next:} Abstract classes, \texttt{final}, and interfaces (design choices)
  \vspace{0.6em}
  \begin{itemize}[<+->]
    \item We will decide when to use abstract class vs interface.
  \end{itemize}
  \slidesources{\OOPSUnitIISources}
\end{frame}

\end{document}
